\documentclass{ctexbook}
\usepackage{amsmath,amssymb,amsthm}
\usepackage[hidelinks]{hyperref}
\newcommand{\corr}[1]{\footnote{本批译校：#1}}
\newcommand{\ednote}[1]{\footnote{英译编者注：#1}}
\newcommand{\Spec}{\operatorname{Spec}}
\newcommand{\Ys}{Y^{\mathrm{sing}}}
\newcommand{\Yss}{Y^{\mathrm{sup\,sing}}}
\title{代数几何基础第五卷预备笔记\\8.7及相邻残段：校订续译}
\author{据用户上传英译编辑稿译出}
\date{2026年11月14日}
\begin{document}
\maketitle
\frontmatter
\chapter*{范围、断裂与处理原则}
本批从底本第28页“Theorem 8.7.”标题起，连续处理该页正文及第29页首段，止于“Corollary 8.9.”标题之前。包括（8.7）(a)、(b)、残缺的(c)、紧随其后的条件列表以及跨页的证明前段。\textbf{不是已经修复并证明的完整（8.7）--（8.8）定理，也未译（8.9）。}

底本为用户上传的 \texttt{EGAVaiello.pdf}，即 Piotr Blass 与 Joseph Blass 翻译编辑的 Grothendieck 第五卷预备笔记。原文(c)中的切空间句子中断，随后直接跳到像的闭包及函数域扩张；页上不见（8.8）标题，但后文（8.9）明确回引（8.8）。不能据此猜造缺失的标题、切空间结论或全套假设。

遵照用户要求：已证实的错误不作为有效数学结论写在正文；原错及理由放脚注。残存而不能定论的文字明确标为“待核残段”，不是把它认证为定理。补充验证与原文翻译分开。对错误等价断言，正文仅列出条件和已核实关系；原等价断言及相应证明文字在脚注中保存，不遗漏其历史内容。

沿用上一批记号：$P=\mathbf P(E)$，$P^\vee=\mathbf P(E^\vee)$，$f:X\to P$ 非分歧，$X/S$ 光滑，$N$ 为法模，$\Ys\simeq\mathbf P(N)$ 通过 $\mathbf P(N(-1))\hookrightarrow\mathbf P(E_X^\vee)$ 嵌入 $Y$。$r=\dim(P/S)$。射影化非空的等维部分上，$\Ys/S$ 的相对维数是 $r-1$。相对微分条件与术语沿用原稿，但“普通奇异点”须与作者后文的术语更正结合阅读。

核查原对话及Library未见本范围已有交付。原对话仍以第七版918页、第0章（19.6.7）为最新结果；独立第五卷分支上一批恰止于（8.7）之前。本批不覆盖已有文件、不自动合并。
\mainmatter
\chapter*{8.7\quad 定理的可读部分}
\noindent\textbf{(a)} 在目前的情形下，子预概形 $\Ys$（定义见第16号或第16节）正是 $\mathbf P(N)$，按前面说明的方式将它视为 $Y$ 的子概形。

\noindent\textbf{(b)} 预概形 $\Yss$ 的底空间（参见第16号或第16节），正是下述态射的分歧点集合：
\[
q:\Ys=\mathbf P(N)\longrightarrow P^\vee.
\]
这里源与靶均相对于 $S$ 光滑，非空的相关部分的相对维数分别为 $r-1$ 和 $r$。也就是说，这个态射在点 $y$ 处非分歧（参见定义），当且仅当 $y$ 对于 $\phi_\xi$ 是几何奇异而非超奇异零点（即本处二次型非退化的奇异零点），其中 $\xi$ 是 $y$ 在 $P^\vee$ 中的像。\corr{原句用的是“geometrically an ordinary singular point”，即“几何普通奇异点”。底本第73--74页作者另行说明：此前把非退化二次奇异零点叫作“ordinary”与惯用术语冲突，应改称“singular non-supersingular zero”或“singular quadratic non-degenerate zero”。故本批采用“非超奇异零点”，限定为这里极化双线性型非退化的二次条件，不把它无条件等同于惯用的普通双点。特征二时奇数个变量的极化型总退化，而 $u^2+vw=0$ 的原点仍是惯用的普通二次奇异点；这一术语差异不能忽略。此处引用后文只为术语校勘，不计作已连续翻译后文。}

\bigskip
\noindent\textbf{(c)〔待核残段〕} 假设 $S=\Spec(k)$，且 $y\in\Ys=\mathbf P(N)$ 为一个 $k$-有理点。令〔字迹不清〕\ednote{请询问 A.G.。〔底本脚注40。〕}和 $\xi$ 分别为它在 $X(k)$ 和 $P^\vee(k)$ 中的投影；考虑 $P^\vee$ 中作为切映射之“像”的线性子簇 $H^1$……

\noindent〔原句在此与下一段内容混接，没有给出可辨认的完整切空间结论。〕\corr{原文从“linear subvariety $H^1$ … image of the tangent map”直接连到“of the closure of its … image in $P^\vee$”，其语法、定义和论证均有断裂；并附法文“son image”（它的像）。这里的省略号明确表示底本连接残缺，不代表主动略去一个已读懂的证明。本批没有补造常见的对偶切空间定理。}

\chapter*{紧随(c)的条件列表与论证残段}
\noindent\textbf{〔编号及假设待核〕} 取其在 $P^\vee$ 中的像的闭包，赋予诱导的约化结构，并考虑诱导态射
\[
g:\Ys\longrightarrow T,
\]
它是整预概形之间的支配态射。\corr{底本此处未完整引入 $T$，也未明确“其”指代的映射。依该句自身及后续记号，$T$ 是所讨论像闭包的名称；但本批不据此认定遗失假设已经恢复。前文 $X/S$ 光滑并不自动保证 $X$ 或 $\Ys$ 整，故必须把“整预概形”视为本段另行要求的条件。后文回指（8.8），本页却没有其标题，本批不擅自补号。}

原稿随后列出下列条件。为清楚起见，$L/K$ 表示由这个支配态射诱导的函数域扩张，$K=k(T)$、$L=k(\Ys)$。\corr{函数域在整性及非空条件下有定义。底本仅说“the field extension $L/K$ defined by $g$”，本批显写其含义，不为其他缺失假设作补充。}
\begin{enumerate}
\item[(i)] 态射 $g$ 泛平展，即至少在一点平展；等价地，在 $\Ys$ 的泛点平展，也就是在该泛点非分歧。
\item[(i bis)] $L/K$ 为有限可分扩张。
\item[(i ter)] $g$ 是双有理态射，即 $L/K$ 为平凡扩张。
\item[(ii)] 从集合论意义说，$\Ys\ne\Yss$。
\item[(ii bis)] 存在 $x\in X(k)$，以及在 $x$ 处与 $X_k$ 相切但在该点不密切的超平面 $H$；这里“不密切”的确切含义是：对定义 $H$ 的 $\mathcal O_{X_k}(1)$ 的截面而言，$x$ 不是超奇异点。
\end{enumerate}

在上述整性、支配性及有限型背景下，可以确认的关系是
\[
\text{(i)}\ \Longleftrightarrow\ \text{(i bis)},
\qquad
\text{(i ter)}\ \Longrightarrow\ \text{(i)}.
\]
原稿要求的全组等价关系，在现存假设下不能成立，须等待缺文核查。\corr{原句译为：“条件(i)与(ii)(bis)是等价的”，但随后实际列出五个条件，证明还要求证明(i)$\Rightarrow$(i ter)。后一蕴涵在仅有“$f$非分歧”时为假。取 $k=\mathbb C$、$X=\mathbb G_m$、$f(t)=[1:t^2:t^4]$。此映射是二次平展覆盖后接到圆锥曲线开子集的同构，再接该子集到 $\mathbf P^2$ 的浸入，故非分歧；$\Ys\simeq X$，而 $g(t)=[t^4:-2t^2:1]$。其像闭包为光滑对偶圆锥曲线，$L/K=\mathbb C(t)/\mathbb C(t^2)$ 次数为二且可分，因此(i)成立而(i ter)不成立。必须补充能排除此类多重覆盖的适当假设；例如 $f$ 为浸入可排除此反例，但本批不恢复缺失的原假设。}

此外，在未假设 $k$ 代数闭的现存文字下，(ii)也不能直接替换为存在 $k$-有理点的(ii bis)。\corr{非空的几何开集未必有 $k$-有理点。例如取 $k=\mathbb R$，$X\subset\mathbf P^2$ 为光滑圆锥曲线 $x_0^2+x_1^2+x_2^2=0$。它没有实点，故(ii bis)失败；但在 $\mathbb C$ 上每条切线的接触均是非退化二重接触，故 $\Yss=\varnothing$ 而 $\Ys\ne\varnothing$，(ii)成立。可通过使用几何点或补入代数闭条件修正这一关系，但原稿缺文未能确定作者究竟采用哪种表述。}

当 $\Ys$ 非空、整、维数 $r-1$，且(ii)成立时，作为分歧闭集的 $\Yss$ 是真闭子集，因而若其非空，便有
\[
\dim\Yss\leq r-2.
\]
所以它在 $P^\vee$ 中的像的闭包余维至少为二。\corr{底本此句印成：“这些条件蕴含 $\Yss\ne\varnothing$〔法文难辨，请询问 A.G.〕$\dim\Yss\leq r-2$，从而其像余维至少二。”非空性不能由前述条件推出；例如在代数闭特征零域上的光滑非退化圆锥曲线没有超奇异切截面。因此正文仅保留经维数论可核实的条件式：$\Yss$ 为空，或其维数至多为 $r-2$。}
若(i bis)成立，还得到
\[
\text{(iii)}\qquad\dim T=r-1,
\]
即 $T$ 在 $P^\vee$ 中余维为一。

\bigskip
\noindent\textbf{证明的现存前段。}
(i)与(i bis)的等价是直接的。至于它与(ii)的等价，需要把 $g$ 到像闭包的非分歧性与 $q:\Ys\to P^\vee$ 的非分歧性在泛点比较；在本段的整性和有限型背景下，可由（8.7）(b)作此比较。\corr{底本原句：“它与(ii)的等价是（8.7）b)的一个直接推论；最后(ii)与(ii bis)的等价，实际上就是 $H^{\mathrm{sup\,sing}}$ 的定义。”后一段中的 $H^{\mathrm{sup\,sing}}$ 疑为 $Y^{\mathrm{sup\,sing}}$，且其有理点断言需要上面说明的基域条件，故不将其作为无条件的正确证明保留在正文。对前一比较，$T\hookrightarrow P^\vee$ 为闭浸入；若 $g$ 泛平展则复合泛非分歧；反之复合泛非分歧蕴含 $g$ 泛非分歧，再由函数域有限可分性得到泛平展。}
显然，(i ter)蕴涵(i)。其逆命题在现存假设下为假，所以不能按原稿所写继续完成一个无条件的证明。\corr{以下完整保存原稿的后续证明文字，而不把它作为有效证明：“还须证明(i)$\Rightarrow$(i ter)。显然可以假设 $K$ 是代数闭的；考虑到假设(i)，归结为证明存在一个非空开集 $U$，使得对每个 $\xi\in U(K)$，在 $\xi$ 上方恰有一个 $\Ys(K)$ 的点。这将由（8.7）c)推出；后者更精确地蕴含下面的断言。”原文将先前函数域 $K=k(T)$ 与基域混用，不能把一个一般函数域无说明地“假设为代数闭”；即便将此理解为对基域作代数闭扩张，二次覆盖反例仍不消失。它所援引的(c)切空间结论在此底本中残缺。因此不能据这段话填补唯一性证明。底本末处(i er)显为(i ter)的排印错，本脚注按所列条件编号书写。}

\backmatter
\chapter*{独立核查记录}
\section*{次数二的反例逐步核验}
令 $s=t^2$。在 $\mathbb C$ 上，$\mathbb G_m\to\mathbb G_m$ 的微分为 $2t\,dt$，处处非零，且该映射有限平展、次数二。映射 $s\mapsto[1:s:s^2]$ 是非退化光滑圆锥曲线的一个开子集的闭嵌入到相应开集。故复合 $f$ 是非分歧态射；$X$ 光滑且整，符合现存背景。

圆锥曲线在 $[1:s:s^2]$ 的切线方程为
\[
s^2X_0-2sX_1+X_2=0.
\]
拉回到参数 $u\in X$，其方程成为
\[
t^4-2t^2u^2+u^4=(u^2-t^2)^2.
\]
在 $u=t$ 和 $u=-t$ 处，它均为非退化二重零点，因为局部二次项系数 $4t^2$ 非零。这也直接说明一个切超平面的非超奇异接触点可以有两个。余法模秩为一，故 $\Ys\simeq X$；参数化的像闭包满足 $b^2=4ac$，为光滑对偶圆锥曲线。映射 $g$ 的函数域扩张为 $\mathbb C(t)/\mathbb C(t^2)$，不是次数一。因此反例不依赖对残缺切空间句子的任何猜测。

\section*{可核实部分与不能声称的部分}
本批(a)承接上批已核实的概形等式；(b)中的“非超奇异”采用作者后文更正后的术语。条件列表完整保留，原错误等价及错误非空断言移至脚注，且提供严格反例。对缺失的(c)结论、缺失的（8.8）标题和假设，不造新原文。

原文最后一段证明在本批已完整处理，但并未因此得到原稿声称的完整定理。特别不能将“已翻译至（8.9）之前”误说成“（8.8）已经严谨证明”。旧批次及全书未在本轮整体回溯修订。

\textbf{下一断点：}底本第29页“Corollary 8.9.”标题。续译前必须沿用本批关于非分歧与浸入的区别、基域与函数域的区别，以及缺失(c)的记录。（8.9）在仅有非分歧假设时亦会受到同一二次覆盖反例的阻碍，不得径直将其开浸入结论当作已核实事实。
\end{document}
